\frontchapter{X. Suggested Priorities for Sample Chapters}
Complete the following six chapters first to test the book's different approaches:

\begin{enumerate}
\item \textbf{Chapter 11, ``What We Say Often Means More Than the Words''}: Test whether school conversations can lead into pragmatics and power.
\item \textbf{Chapter 19, ``Who Speaks, and Who Sees?''}: Test whether close literary reading can move beyond prescribed answers.
\item \textbf{Chapter 28, ``Why Did Humans Begin Farming?''}: Test whether global history can resist simplistic progress narratives.
\item \textbf{Chapter 46, ``Studying Religion Without Rushing to Judgment''}: Test respect, comparison, and the limits of evidence in sensitive subjects.
\item \textbf{Chapter 63, ``What Makes an Action Right?''}: Test whether philosophy can move from everyday dilemmas into serious theory.
\item \textbf{Chapter 96, ``Is Beauty in the Object or in the Beholder?''}: Test whether aesthetic experience and argument can be combined.
\end{enumerate}
Suggested sample length: 8,000--12,000 Chinese characters. Invite students in the second and third years of middle school, language or history teachers, and adults without an academic humanities background to read them. Record especially:

\begin{itemize}
\item Which concept first makes readers lose the picture;
\item Which story entertains without advancing the question;
\item Which theoretical passage feels like memorizing terminology;
\item Which comparison erases differences between traditions;
\item Which value judgment is mistaken for a historical fact;
\item Which open question makes readers want to keep discussing.
\end{itemize}
